\documentclass[11pt,reqno]{amsart}
\usepackage[localtheorems]{raeez-math-template}

% ================================================================
% Packages
% ================================================================
\usepackage{array}

% ================================================================
% Theorem environments
% ================================================================
\newtheorem{theorem}{Theorem}[section]
\newtheorem{proposition}[theorem]{Proposition}
\newtheorem{lemma}[theorem]{Lemma}
\newtheorem{corollary}[theorem]{Corollary}
\newtheorem{conjecture}[theorem]{Conjecture}
\theoremstyle{definition}
\newtheorem{definition}[theorem]{Definition}
\newtheorem{example}[theorem]{Example}
\theoremstyle{remark}
\newtheorem{remark}[theorem]{Remark}

% ================================================================
% Macros
% ================================================================
\newcommand{\cA}{\mathcal{A}}
\newcommand{\cC}{\mathcal{C}}
\newcommand{\cM}{\mathcal{M}}
\newcommand{\cZ}{\mathcal{Z}}
\newcommand{\cD}{\mathcal{D}}
\newcommand{\cO}{\mathcal{O}}
\newcommand{\cW}{\mathcal{W}}
\newcommand{\barB}{\bar{B}}
\newcommand{\barBgeom}{\bar{B}^{\mathrm{ch}}}
\newcommand{\Ran}{\mathrm{Ran}}
\newcommand{\MC}{\mathrm{MC}}
\newcommand{\Sym}{\mathrm{Sym}}
\newcommand{\Symch}{\mathrm{Sym}^{\mathrm{ch}}}
\newcommand{\Hom}{\mathrm{Hom}}
\newcommand{\Ext}{\mathrm{Ext}}
\newcommand{\Res}{\mathrm{Res}}
\newcommand{\FM}{\overline{C}}
\newcommand{\fg}{\mathfrak{g}}
\newcommand{\fh}{\mathfrak{h}}
\newcommand{\Vir}{\mathrm{Vir}}
\newcommand{\KM}{\mathrm{KM}}
\newcommand{\Heis}{\mathrm{Heis}}
\newcommand{\ChirHoch}{\mathrm{ChirHoch}}

\DeclareMathOperator{\gr}{gr}

\numberwithin{equation}{section}

% ================================================================
\begin{document}

\title[Koszulness chart torsors]
{Koszulness chart torsors and chiral Koszulness criteria}

\author{Raeez Lorgat}
\address{Perimeter Institute for Theoretical Physics,
 Waterloo, ON N2L 2Y5, Canada}
\email{rlorgat@perimeterinstitute.ca}

\date{}

\begin{abstract}
For a chiral algebra $\cA$ on a smooth projective curve~$X$ with
PBW filtration, the Koszulness locus is a derived formal moduli
problem over the Drinfeld associator torsor. At finite PBW stage
and finite $\mathrm{GRT}_1$ weight quotient, after fixing an
associator coordinate and a gauge framing, its classical truncation
is an open subscheme of a finite-type affine Maurer--Cartan scheme.
Unframed finite stages are quotient stacks; the full completed
object is a derived inverse limit. Nonempty framed components carry
the resulting pro-$\mathrm{GRT}_1$ chart torsor.

The intrinsic Koszulness core consists of the chiral Koszul
twisting datum, PBW $E_2$ collapse, $A_\infty$ formality of the
bar model, and bar--cobar counit inversion in the Theorem-B
ambient. Ext diagonal concentration is a consequence, with a
converse under the chiral BGS/Priddy diagonal-detection
hypothesis. Kac--Shapovalov nondegeneracy is a family criterion.
Barr--Beck--Lurie monadicity is a consequence of the counit.
Factorization homology is a criterion only under the stated
comparison and detection hypotheses. Chiral Hochschild
concentration is a one-way consequence on the Koszul locus; a
free-polynomial cup-product description requires Massey vanishing.
The Lagrangian criterion is conditional on perfectness and
nondegeneracy, and $\cD$-module purity is used in the proved
direction toward boundary acyclicity.
\end{abstract}

\maketitle

% ================================================================
\section{Introduction}
% ================================================================

\subsection{The problem}

The Koszul duality of a graded associative algebra
$A = T(V)/(R)$ with quadratic relations $R \subseteq V \otimes V$
admits many equivalent formulations.
Priddy's original criterion is purity of the
bar complex~\cite{Priddy}. Beilinson, Ginzburg, and Soergel extend
this to the derived category of a highest-weight category and give
a cohomological characterization via Ext diagonal
vanishing~\cite{BGS}. Loday and Vallette place Koszul duality in
the operadic framework and reformulate it as an
$A_\infty$-formality condition for the Koszul dual
cooperad~\cite{LV12}. Polishchuk and Positselski develop the dual
story for corings and comodules~\cite{PP}. Each of these
reformulations, in the associative world, has a parallel in the
world of chiral algebras (equivalently, vertex algebras) over an
algebraic curve: the bar complex carries an extra layer of
factorization geometry, the Koszul dual inherits an
$\cO$-module structure, and the relevant spectral sequence lives
on a Fulton--MacPherson compactification.

The classical web of Koszul criteria survives in the chiral world
only after its scope is separated. Seven criteria form the
intrinsic bidirectional core. Kac--Shapovalov nondegeneracy and
Fulton--MacPherson boundary acyclicity are chiral criteria inside
that core. Barr--Beck--Lurie monadicity is a consequence of the
bar--cobar counit. Factorization homology requires comparison and
detection hypotheses. Chiral Hochschild concentration is a
one-way theorem on the Koszul locus. The Lagrangian and
$\cD$-module tests retain their own perfectness and purity
hypotheses.

\subsection{Koszulness moduli}

The Koszulness moduli object is controlled by the completed
convolution dg Lie algebra of the bar--cobar counit. It is a
derived formal moduli problem before truncation. A scheme appears
only after finite PBW truncation, finite $\mathrm{GRT}_1$ weight
quotient, associator coordinate, and gauge framing.

\begin{definition}[Koszulness moduli object]
\label{def:koszulness-moduli-N1}
For a complete local Artinian commutative dg $\bQ$-algebra~$R$
with maximal ideal~$\mathfrak m_R$, set
\[
\widehat{\mathfrak g}_{\Phi}(\cA)
=
\varprojlim_N
\mathrm{Conv}_{\Phi}\bigl(F_{\leq N}\barBgeom(\cA),
F_{\leq N}\cA\bigr).
\]
The Koszulness moduli object is the derived functor
\[
\cM_{\mathrm{Kosz}}(\cA)\colon
\mathbf{Art}^{\wedge}_{\bQ}\longrightarrow\mathbf{Spaces},
\]
\[
R\longmapsto
\operatorname{hofib}_{0}
\left[
\coprod_{\Phi\in\mathrm{Assoc}(R)}
\mathrm{MC}_{\bullet}\bigl(
\widehat{\mathfrak g}_{\Phi}(\cA)\widehat{\otimes}\mathfrak m_R
\bigr)
\xrightarrow{\;\xi\mapsto\operatorname{Cone}(\varepsilon_\xi)\;}
D^{\mathrm{co}}_{\mathrm{ch}}(X;R)
\right],
\]
where
$\varepsilon_\xi\colon
\Omega^{\mathrm{ch}}_{\Phi}(\widehat{\barBgeom}(\cA),\xi)
\to \cA\widehat{\otimes}R$
is the deformed counit. An $R$-point is a triple
$(\Phi,\xi,h_\xi)$ consisting of an associator, a Maurer--Cartan
element, and a trivialisation of
$\operatorname{Cone}(\varepsilon_\xi)$ in the coderived chiral
category over~$R$.
\end{definition}

\begin{theorem}[Koszulness moduli, finite-stage representability]
\label{thm:meta-as-moduli-N1}
Let $\cA$ be of finite PBW type. At each finite PBW stage and
finite weight quotient of $\mathrm{GRT}_1$, assume PBW--Artin
type: the finite convolution dg Lie algebra is nilpotent with
finite-dimensional bounded cohomology, the gauge group is
unipotent algebraic, the finite Koszul cone is perfect with
coherent cohomology, and the tower is derived Mittag--Leffler.
Then $\cM_{\mathrm{Kosz}}(\cA)$ is a possibly empty derived formal
moduli problem over the Drinfeld associator torsor. The unframed
finite-stage truncation is a derived Artin quotient stack. After
fixing an associator coordinate and a gauge framing, its classical
truncation is an open subscheme of a finite-type affine
Maurer--Cartan scheme, cut out by acyclicity of the finite Koszul
cone. The completed object is the derived inverse limit of these
finite stages.

The moduli object is nonempty if and only if the completed
bar--cobar counit
$\Omega^{\mathrm{ch}}(\widehat{\barBgeom}(\cA))\to\cA$
is a quasi-isomorphism. On each nonempty framed component, the
projection to the associator torsor is a pro-$\mathrm{GRT}_1$
torsor. Before framing the same component is the quotient of that
torsor by the automorphism group of the witness.
\end{theorem}

\subsection{Classical Koszulness criteria}

Let $\cA$ be a chiral algebra on a smooth projective curve~$X$,
equipped with its Poincar\'e--Birkhoff--Witt filtration
$F_\bullet\cA$.

\begin{theorem}[Equivalences and consequences of chiral Koszulness]
The following intrinsic conditions are equivalent:
\begin{enumerate}[leftmargin=2.8em]
\item[\textup{(i)}] $\cA$ is chirally Koszul in the sense of
 Definition~\ref{def:chiral-koszul-morphism};
\item[\textup{(ii)}] the PBW spectral sequence on $\barBgeom(\cA)$ collapses
 at~$E_2$;
\item[\textup{(iii)}] the minimal $A_\infty$-model of $\barBgeom(\cA)$ is formal:
 $m_n = 0$ for $n \ge 3$;
\item[\textup{(iv)}] $\Ext^{p,q}_\cA(\omega_X,\omega_X) = 0$ for $p \ne q$;
\item[\textup{(v)}] in the Theorem-B ambient, with
 $C_\cA=\barBgeom(\cA)$ on raw finite PBW windows and
 $C_\cA=\widehat{\barBgeom}(\cA)$ on the strict completed
 pro-conilpotent regime, the continuous bar--cobar counit
 $\Omega_X^{\mathrm{cont}}(C_\cA) \to \cA$ is a quasi-isomorphism;
\item[\textup{(ix)}] the Kac--Shapovalov determinant $\det G_h \neq 0$ in the
 bar-relevant range;
\item[\textup{(x)}] for every Fulton--MacPherson boundary stratum~$S$, the
 restriction $i_S^!\,\barB_n(\cA)$ is acyclic outside degree~$0$.
\end{enumerate}
The Barr--Beck--Lurie comparison functor for the bar--cobar
adjunction is a consequence of the counit criterion. Factorization
homology concentration is a criterion only under the comparison and
detection hypotheses stated in Theorem~\ref{thm:koszul-equivalences-meta}.
On the Koszul locus, $\ChirHoch^*(\cA)$ has cohomological amplitude
$[0,2]$ and Hochschild duality; a free-polynomial cup-product
description is conditional on Massey vanishing.
Under a perfectness and non-degeneracy hypothesis on the ambient
tangent complex, condition
\begin{enumerate}[label=\textup{(\roman*)},leftmargin=2.4em,start=11]
\item the Koszul locus subfunctor $\cM_\cA$ and its dual
 $\cM_{\cA^!}$ are transverse Lagrangians in the $(-1)$-shifted
 symplectic deformation space $\cM_{\mathrm{comp}}$
\end{enumerate}
is equivalent to them. Finally,
\begin{enumerate}[label=\textup{(\roman*)},leftmargin=2.4em,start=12]
\item each $\barBgeom_n(\cA)$ is pure of weight~$n$ as a mixed
 Hodge module, with characteristic variety aligned to the FM
 boundary strata
\end{enumerate}
implies \textup{(x)}; the converse is proved for the affine
Kac--Moody family, and open in general.
\end{theorem}

The affine Kac--Moody family is controlled by the Kac--Shapovalov
determinant. The $\beta\gamma$ system is controlled by the free
PBW bar model. The Virasoro Koszul locus is controlled by PBW
collapse. These three tests land in the same intrinsic core.

\subsection{History and related work}

Priddy~\cite{Priddy} introduced the Koszul property for quadratic
algebras via resolutions. Beilinson, Ginzburg, and
Soergel~\cite{BGS} showed that Koszulness is equivalent to Ext
diagonal vanishing in a broad categorical context, a formulation
that survives in (iv). Polishchuk and Positselski~\cite{PP}
carried the theory to noncommutative, graded, and curved settings.
Loday and Vallette~\cite{LV12} organized the operadic framework
and characterized Koszulness as $A_\infty$-formality of the Koszul
dual cooperad, matching condition~\textup{(iii)}. Positselski's construction
of the Koszul duality functors between modules and comodules is
the associative prototype of the chiral bar--cobar functors used
below.

On the chiral side, Beilinson and Drinfeld~\cite{BD04} developed
factorization algebras and chiral algebras on a curve and isolated
the role of the PBW filtration. Francis and
Gaitsgory~\cite{FG12} gave an adjunction between chiral Lie and
chiral commutative operads on $\Ran(X)$, which is the operadic
backbone of the bar--cobar adjunction here. Gaitsgory and
Rozenblyum~\cite{GR17} give the $\infty$-categorical framework
adequate to formulate condition~(vi). Costello and
Gwilliam~\cite{CG} build the factorization homology machinery
used in~(vii). On the Kac--Shapovalov side, the vanishing of the
Shapovalov determinant in specific ranges is a classical criterion
for irreducibility~\cite{KacBook}; its role in controlling the
PBW spectral sequence was hinted at in the work of
Feigin--Fuchs on Virasoro representations.

The synthesis is the derived chart atlas
$\cM_{\mathrm{Kosz}}(\cA)$ together with the intrinsic equivalence
core of Theorem~\ref{thm:koszul-equivalences-meta}.

\subsection{Conventions}

We work in characteristic zero. All chain complexes are
cohomologically graded, $|d| = +1$. Desuspension lowers degree:
$|s^{-1}v| = |v| - 1$. The bar complex uses the augmentation
ideal $\bar\cA = \ker(\varepsilon)$, not $\cA$ itself; we write
$\barBgeom(\cA) = T^c(s^{-1}\bar\cA)$ for the geometric chiral bar
complex, with deconcatenation coproduct. When we refer to the
symmetric bar $\barB^\Sigma(\cA)$, we specify it explicitly.
Imported results from the monograph are cited at the point where
their hypotheses enter.

% ================================================================
\section{Preliminaries}
\label{sec:preliminaries}
% ================================================================

\subsection{Chiral algebras and the PBW filtration}

Let $X$ be a smooth projective curve over an algebraically closed
field of characteristic zero. A \emph{chiral algebra} on~$X$ is a
left $\cD_X$-module $\cA$ equipped with a chiral bracket
\[
\mu\colon j_*j^*(\cA\boxtimes\cA) \longrightarrow
 \Delta_*\cA,
\]
where $j\colon X^2\setminus\Delta\hookrightarrow X^2$ is the
complement of the diagonal and $\Delta\colon X\to X^2$ is the
diagonal embedding. The bracket is required to satisfy the
Beilinson--Drinfeld axioms of skew-symmetry and Jacobi, equivalent
in the vertex algebra language to the Borcherds
identity~\cite{BD04}.

We write $\varepsilon\colon \cA \to \omega_X$ for the augmentation
(the projection to the vacuum), and
$\bar\cA = \ker(\varepsilon)$ for the augmentation ideal.
The \emph{PBW filtration} $F_\bullet\cA$ is the increasing
filtration generated by placing the generators in $F_1\cA$ and
closing under the chiral bracket. We write
$R_\cA = \gr^F_\bullet\cA$ for the associated graded; it is a
Poisson vertex algebra~\cite{Li}. The filtration is bounded below
(by $F_0 = \omega_X$) and exhaustive.

\subsection{The geometric bar complex}

Fix $n \ge 1$. The bar complex in degree~$n$ is constructed on the
Fulton--MacPherson compactification $\FM_n(X)$ of the
configuration space $C_n(X) = X^n\setminus\bigcup_{i<j}\Delta_{ij}$.
As a graded $\cD_{\FM_n(X)}$-module,
\[
\barBgeom_n(\cA)
\;=\; j_*j^*\bar\cA^{\boxtimes n}
 \otimes \Omega^{n-1}_{\FM_n(X)}(\log D)[1-n],
\]
where $D$ is the boundary divisor and the shift $[1-n]$ is the
desuspension grading $|s^{-1}v| = |v|-1$ applied $n-1$ times.
The assembly
\[
\barBgeom(\cA) \;=\; \bigoplus_{n \ge 1} \barBgeom_n(\cA)
\]
carries the deconcatenation coproduct and a differential
$d_{\mathrm{bar}} = d_\cA + d_{\mathrm{OPE}} + d_{\mathrm{dR}}$
whose three components encode the internal $\cD$-differential, the
OPE residue along the collision divisors, and the de~Rham
differential on $\FM_n(X)$~\cite{BD04,LorgatMKD1}.

\begin{remark}
The use of the augmentation ideal $\bar\cA$ is essential: writing
$T^c(s^{-1}\cA)$ in place of $T^c(s^{-1}\bar\cA)$ would include the
vacuum and destroy the bar-cobar adjunction
because the vacuum summand is not part of the reduced bar coalgebra.
\end{remark}

\subsection{The cobar functor}

Dually, for a factorization coalgebra $\cC$ on~$X$ with
coaugmentation $\eta\colon \omega_X\to\cC$, the cobar complex is
\[
\Omega^{\mathrm{ch}}(\cC)
\;=\; \bigoplus_{n \ge 1}
 (s\bar\cC)^{\otimes n}\otimes \Omega^{n-1}_{\FM_n(X)}(\log D)[n-1],
\]
where $\bar\cC = \mathrm{coker}(\eta)$ and the differential is the
dual of the bar differential.

\begin{proposition}[Bar--cobar adjunction;~\cite{LorgatMKD1}]
\label{prop:bar-cobar-adjunction}
The functors $\barBgeom\dashv\Omega^{\mathrm{ch}}$ form an adjoint
pair between chiral algebras and factorization coalgebras on~$X$.
The unit
$\eta_\cC\colon \cC\to \barBgeom(\Omega^{\mathrm{ch}}(\cC))$
and counit
$\varepsilon_\cA\colon
 \Omega^{\mathrm{ch}}(\barBgeom(\cA))\to\cA$
are natural transformations.
\end{proposition}

\subsection{Chiral twisting morphisms and Koszulness}

Let $\tau\colon \cC\to \cA$ be a twisting morphism: a degree~$+1$
map of $\cD_X$-modules satisfying the Maurer--Cartan equation
$d\tau + \tau\star\tau = 0$ with respect to the convolution
product. The \emph{left twisted tensor product}
$K_\tau^L(\cA,\cC)$ is the complex
$\cA\otimes\cC$ with differential
$d_\cA\otimes 1 + 1\otimes d_\cC + (\mu\otimes 1)(1\otimes\tau\otimes 1)
 (1\otimes \Delta_\cC)$;
the right version $K_\tau^R(\cC,\cA)$ is defined symmetrically.

\begin{definition}[Chiral Koszul morphism]
\label{def:chiral-koszul-morphism}
A chiral twisting datum $(\cA,\cC,\tau,F_\bullet)$ is
\emph{Koszul} if:
\begin{enumerate}[label=\textup{(\alph*)}]
\item $\gr_F\cA$ is classically Koszul on each finite PBW window;
\item both $K_\tau^L(\cA,\cC)$ and $K_\tau^R(\cC,\cA)$
are acyclic;
\item the PBW filtration is complete, separated, and
Mittag--Leffler on finite windows, so associated-graded
acyclicity lifts to the completed chiral bar complex.
\end{enumerate}
The chiral algebra~$\cA$ is \emph{chirally Koszul}
if there exists a Koszul twisting datum of the form
$(\cA,\barBgeom(\cA),\tau_0,F_\bullet^{\mathrm{PBW}})$ where
$\tau_0$ is the canonical bar twisting morphism.
\end{definition}

\subsection{The PBW spectral sequence}

The PBW filtration on $\cA$ induces a filtration on the bar
complex $\barBgeom(\cA)$ by bar degree, producing a spectral
sequence
\[
E_1^{p,q}(\cA) \;=\; H^{p+q}(\gr^F_p \barBgeom(\cA))
 \;\Longrightarrow\; H^{p+q}(\barBgeom(\cA)).
\]
On the $E_1$ page, $\gr^F_p\barBgeom(\cA)$ is computed from the
Poisson vertex algebra $R_\cA$ by the Chevalley--Eilenberg-style
bar construction; the higher differentials $d_r$ ($r \ge 2$)
encode quantum corrections. Collapse at $E_2$ means $d_r = 0$ for
all $r \ge 2$, which is condition~(ii) of
Theorem~\ref{thm:koszul-equivalences-meta}.

\subsection{The chiral Hochschild complex}

The chiral Hochschild complex
$C^\bullet_{\mathrm{ch}}(\cA)$ has degree-$n$ cochains
\[
C^n_{\mathrm{ch}}(\cA) \;=\; \Gamma\bigl(\FM_{n+2}(X),
 j_*j^*\cA^{\boxtimes(n+2)}\otimes
 \Omega^n_{\FM_{n+2}(X)}(\log D)\bigr),
\]
with differential
$d_{\mathrm{Hoch}}
 = d_{\mathrm{int}} + d_{\mathrm{fact}} + d_{\mathrm{config}}$
(internal, factorization, configuration-space de~Rham).
Its cohomology $\ChirHoch^*(\cA)$ is the tangent complex to the
moduli of chiral algebra structures at~$\cA$~\cite{LorgatMKD1}.

% ================================================================
\section{The Koszulness Criteria}
\label{sec:meta}
% ================================================================

\begin{theorem}[Equivalences and consequences of chiral Koszulness]
\label{thm:koszul-equivalences-meta}
Let $\cA$ be a chiral algebra on a smooth projective curve~$X$
with PBW filtration $F_\bullet$.
For a smooth projective curve~$\Sigma$, let
$\mathsf{FH}_{\Sigma}(\cA)$ denote the filtered comparison
between the geometric bar object and the factorization homology
realization, with strong convergence for the bar-length and
Goodwillie filtrations. Let
$\mathsf{FH}^{\mathrm{det}}_{\Sigma}(\cA)$ denote the detecting
strengthening: degree-zero concentration after realization reflects
degree-zero concentration of the geometric bar model.

Conditions \textup{(i)--(iii)} and \textup{(v)} below are the
intrinsic bidirectional equivalence core in the Theorem-B ambient.
Condition~\textup{(iv)} is a consequence of the core; its converse
requires the chiral BGS/Priddy diagonal-detection hypothesis.
Condition~\textup{(vi)} is a consequence of~\textup{(v)}.
Condition~\textup{(vii)} is a factorization-homology consequence
under $\mathsf{FH}$, and becomes a criterion only under
$\mathsf{FH}^{\mathrm{det}}$. Condition \textup{(viii)} is a
one-way chiral Hochschild consequence on the Koszul locus.
Condition~\textup{(ix)} is a bidirectional criterion on the
families where a Kac--Shapovalov determinant controls the PBW
comparison. Condition~\textup{(x)} is equivalent to the core when
FM boundary detection holds. Under additional perfectness and
non-degeneracy hypotheses on the ambient tangent complex,
condition~\textup{(xi)} is also equivalent to the intrinsic core.
Condition~\textup{(xii)} implies condition~\textup{(x)}; the
converse is known for affine Kac--Moody at non-critical level and
is open for Virasoro and $\cW$-algebras.

\smallskip
\noindent\textbf{Intrinsic equivalence core and monadic consequence:}
\begin{enumerate}[leftmargin=2.8em]
\item[\textup{(i)}] $\cA$ is chirally Koszul
 (Definition~\ref{def:chiral-koszul-morphism}).
\item[\textup{(ii)}] The PBW spectral sequence on $\barBgeom(\cA)$ collapses
 at~$E_2$.
\item[\textup{(iii)}] The minimal $A_\infty$-model of $\barBgeom(\cA)$ is formal:
 $m_n = 0$ for $n \ge 3$.
\item[\textup{(iv)}] Ext diagonal vanishing:
 $\Ext^{p,q}_\cA(\omega_X,\omega_X) = 0$ for $p \ne q$.
\item[\textup{(v)}] The bar--cobar counit
 $\Omega_X^{\mathrm{cont}}(C_\cA) \to \cA$ is a
 quasi-isomorphism, with $C_\cA=\barBgeom(\cA)$ on raw finite PBW
 windows and $C_\cA=\widehat{\barBgeom}(\cA)$ on the strict
 completed pro-conilpotent regime.
\item[\textup{(vi)}] The Barr--Beck--Lurie comparison for the bar--cobar
 adjunction is an equivalence on the fiber
 over $\barBgeom(\cA)$.
\item[\textup{(ix)}] The Kac--Shapovalov determinant $\det G_h \neq 0$ in the
 bar-relevant range.
\item[\textup{(x)}] The restriction $i_S^!\,\barB_n(\cA)$ to every FM boundary
 stratum~$S$ is acyclic outside degree~$0$.
\end{enumerate}

\smallskip
\noindent\textbf{Factorization homology and Hochschild consequences:}
\begin{enumerate}[label=\textup{(\roman*)},leftmargin=2.4em,start=7]
\item If $\mathsf{FH}_{\Sigma}(\cA)$ holds, chiral Koszulness
 implies
 \[
 H^k\!\left(\textstyle\int_{\Sigma}\cA\right)=0
 \qquad (k\neq 0).
 \]
 If $\mathsf{FH}^{\mathrm{det}}_{\mathbb P^1}(\cA)$ also holds,
 this genus-zero concentration is equivalent to the intrinsic core.
 Under the uniform-weight hypothesis, the all-genera scalar obstruction is
 $\mathrm{obs}_g(\cA)=\kappa(\cA)\lambda_g$.
\item On the Koszul locus, $\ChirHoch^\bullet(\cA)$ has
 cohomological amplitude $[0,2]$, Hochschild duality with
 $\cA^!$, and Hilbert series determined by the degree
 $0,1,2$ terms. A free-polynomial cup-product description requires
 the corresponding Massey-vanishing hypothesis.
\end{enumerate}

\smallskip
\noindent\textbf{Conditional and partial:}
\begin{enumerate}[label=\textup{(\roman*)},leftmargin=2.4em,start=11]
\item \emph{Lagrangian criterion:} $\cM_\cA$ and $\cM_{\cA^!}$
 are transverse Lagrangians in the $(-1)$-shifted symplectic
 deformation space~$\cM_{\mathrm{comp}}$.
\item $\cD$-\emph{module purity:} each $\barBgeom_n(\cA)$ is pure
 of weight~$n$ as a mixed Hodge module, with characteristic
 variety aligned to FM boundary strata.
\end{enumerate}
\end{theorem}

% ================================================================
\section{Proof of the Koszulness criteria}
\label{sec:proof}
% ================================================================

The intrinsic circuit is
\[
(\mathrm{i}) \Leftrightarrow (\mathrm{ii})
\Leftrightarrow (\mathrm{iii})
\Leftrightarrow (\mathrm{v});
\]
conditions \textup{(iv)}, \textup{(ix)}, and \textup{(x)} are
equivalent to nodes of this circuit. Conditions \textup{(vi)},
\textup{(vii)}, and \textup{(viii)} are consequences with the
hypotheses stated in Theorem~\ref{thm:koszul-equivalences-meta}.
Condition~\textup{(xi)} is treated under the perfectness and
non-degeneracy hypotheses. Condition~\textup{(xii)} is used in the
proved direction toward boundary acyclicity.

\subsection{The core circuit}

\subsubsection*{\textup{(i)}~$\Leftrightarrow$~\textup{(ii)}}

Koszulness, in the sense of
Definition~\ref{def:chiral-koszul-morphism}, is by definition
acyclicity of $K_\tau^L(\cA,\barBgeom(\cA))$ and
$K_\tau^R(\barBgeom(\cA),\cA)$ for the canonical twisting morphism
$\tau_0$. The left twisted tensor product is filtered by the PBW
bar degree, and on the associated graded the twisted tensor product
becomes the standard Koszul complex of the Poisson vertex algebra
$R_\cA$. Acyclicity on the associated graded lifts to acyclicity on
$K_\tau^L$ if and only if the bar spectral sequence collapses at
$E_2$~\cite[Thm.~4.2.1]{LorgatMKD1}. This is the PBW Koszulness
criterion, and it gives the first equivalence.

\subsubsection*{\textup{(ii)}~$\Rightarrow$~\textup{(v)}}

$E_2$-collapse gives bar concentration on the diagonal, which is
the input for bar--cobar inversion. Specifically, if the spectral
sequence collapses, then the cohomology $H^*(\barBgeom(\cA))$ is
concentrated in bidegree $(p,p)$; the cobar functor on a
diagonally concentrated bigraded coalgebra returns the algebra
back, giving
the continuous counit
$\Omega_X^{\mathrm{cont}}(C_\cA)\xrightarrow{\sim}\cA$ in the
Theorem-B ambient. Here $C_\cA=\barBgeom(\cA)$ on raw finite PBW
windows and $C_\cA=\widehat{\barBgeom}(\cA)$ on the strict
completed pro-conilpotent regime.

\subsubsection*{\textup{(v)}~$\Rightarrow$~\textup{(i)}}

If the counit is a quasi-isomorphism, then its mapping cone is
acyclic. The mapping cone of $\varepsilon_\cA$ is canonically
quasi-isomorphic to the twisted tensor product $K_\tau^L$
(a standard lemma about bar--cobar; see~\cite[\S2.2]{LV12} in the
associative case, adapted to the chiral setting
in~\cite{LorgatMKD1}). Acyclicity of $K_\tau^L$ is the left
twisted tensor part of Definition~\ref{def:chiral-koszul-morphism};
the right twisted tensor product follows by the same universal
twisting-morphism theorem and the PBW convergence hypothesis.

\subsubsection*{\textup{(v)}~$\Rightarrow$~\textup{(viii)}}

The quasi-isomorphism $\Omega^{\mathrm{ch}}(\barBgeom(\cA))
\xrightarrow{\sim}\cA$ is a free resolution of~$\cA$ by a chiral
cofree algebra. Computing chiral Hochschild cohomology through
this resolution turns $\ChirHoch^*(\cA)$ into a bigraded complex
whose $E_1$-page is the cotangent complex of the bar coalgebra.
The bar-cobar filtration gives cohomological amplitude $[0,2]$
and the Hochschild duality pairing with the Verdier dual
Koszul algebra~$\cA^!$. A free-polynomial product statement is
stronger; it follows only under the Massey-vanishing hypothesis
for the transferred Fulton--MacPherson product. Chiral Hochschild
concentration alone does not imply the counit criterion here.

\subsubsection*{\textup{(ii)}~$\Leftrightarrow$~\textup{(iii)}}

For the forward direction, $E_2$-collapse of the PBW spectral
sequence means all differentials $d_r = 0$ for $r \ge 2$. The
Homological Perturbation Lemma identifies the differentials $d_r$
of the spectral sequence with the $A_\infty$ operations
$m_{r+1}$ of the minimal model under the transfer along the
PBW filtration, up to sign. Hence $d_r = 0$ for $r \ge 2$ is
equivalent to $m_n = 0$ for $n \ge 3$.

For the reverse direction, if the minimal $A_\infty$-model is
formal then $H^*(\barBgeom(\cA))$ is a strict dg algebra with no
higher operations, and its bar spectral sequence collapses at
$E_2$ by Keller's theorem on formal
$A_\infty$-algebras~\cite{Keller}. The argument proceeds
fiberwise on FM strata; each fiber is an $A_\infty$-algebra over
a field, to which Keller's theorem applies directly. The PBW
filtration is compatible with the FM stratification, so fiberwise
collapse assembles to global collapse.

\subsection{Ext diagonal vanishing and the core}

\subsubsection*{\textup{(i)}~$\Rightarrow$~\textup{(iv)}}

$E_2$-collapse places $H^*(\barBgeom(\cA))$ on the diagonal
$p = q$; the bar complex computes
$\Ext^{p,q}_\cA(\omega_X,\omega_X)$, so
$\Ext^{p,q}_\cA(\omega_X,\omega_X) = 0$ for $p \neq q$.

\subsubsection*{\textup{(iv)}~$\Rightarrow$~\textup{(i)}}

Assume the chiral BGS/Priddy diagonal-detection hypothesis for
the chosen finite PBW component. If the Ext bigrading is diagonal,
then
$\barBgeom(\cA)$ has cohomology concentrated on the line $p = q$.
A $d_r$-differential in the PBW spectral sequence moves
off-diagonal by $(r, 1-r)$, which would produce off-diagonal
classes, contradicting diagonal concentration of the abutment.
Hence $d_r = 0$ for all $r \ge 2$.

The spectral sequence converges conditionally in the sense of
Boardman (filtration Hausdorff and complete, since the PBW
filtration on a positively graded chiral algebra is bounded
below). In characteristic zero with a split filtration (the PBW
filtration splits by conformal weight), the extension problem is
trivial: $E_\infty = E_2$ as bigraded objects. So $E_2$-collapse
gives~(ii), which is equivalent to~(i).

\subsection{Factorization homology concentration}

\subsubsection*{\textup{(i)}~$\Rightarrow$~\textup{(vii)}}

Assume $\mathsf{FH}_{\Sigma}(\cA)$. The filtered comparison
identifies the associated graded of the Goodwillie filtration on
$\int_{\Sigma}\cA$ with the corresponding geometric bar
contributions. PBW collapse concentrates those contributions in
degree~$0$, hence
$H^k(\int_{\Sigma}\cA)=0$ for $k\neq 0$. Under the
uniform-weight hypothesis, the surviving scalar obstruction is
$\mathrm{obs}_g(\cA)=\kappa(\cA)\lambda_g$.

\subsubsection*{\textup{(vii)}~$\Rightarrow$~\textup{(i)} under detection}

Assume in addition
$\mathsf{FH}^{\mathrm{det}}_{\mathbb P^1}(\cA)$. Then
degree-zero concentration of $\int_{\mathbb P^1}\cA$ reflects
degree-zero concentration of $\barBgeom_{\mathbb P^1}(\cA)$.
This is PBW $E_2$ collapse, hence condition~\textup{(ii)} and
therefore chiral Koszulness.

\subsection{FM boundary acyclicity}

\subsubsection*{\textup{(i)}~$\Rightarrow$~\textup{(x)}}

Each FM boundary stratum
$S_T \cong \prod_{v\in V(T)}\FM_{|v|}(X_v)$ is indexed by a
rooted tree~$T$. The residue component $d_{\mathrm{res}}|_{S_T}$
of the bar differential restricts to the tensor product of
lower-degree bar complexes, one for each vertex of~$T$.
$E_2$-collapse at each vertex degree gives acyclicity at $S_T$:
$H^k(i_{S_T}^!\,\barB_n(\cA)) = 0$ for $k \neq 0$.

\subsubsection*{\textup{(x)}~$\Rightarrow$~\textup{(i)}}

At the deepest stratum, the binary collision divisor
$S = D_{\{1,2\}}\cong X\times\FM_{n-1}(X)$, the restriction
$i_S^!\,\barB_n(\cA)$ computes the bar-complex contribution from
the OPE $a_{(k)}b$ at a single collision point. Acyclicity of
$i_S^!$ at every stratum forces every residue to be exact, which
is the PBW condition $d_1(e_i) = 0$ on the associated graded.
Iterating over all strata gives $E_2$-collapse, hence~(i).

\subsection{Barr--Beck--Lurie monadicity}

\subsubsection*{\textup{(v)}~$\Rightarrow$~\textup{(vi)}}

The Barr--Beck--Lurie theorem~\cite[Thm.~4.7.3.5]{HA} states that
the comparison functor $\Phi$ associated with the bar--cobar
adjunction is an equivalence if and only if the bar functor
$\barB_\kappa$ is conservative and preserves totalizations of
$\barB_\kappa$-split cosimplicial objects.

Conservativity holds on the Koszul locus: $\barB_\kappa$ detects
quasi-isomorphisms, because the bar--cobar counit inverts them;
this is condition~\textup{(v)}.

Totalization preservation: at genus~$0$ this follows from the
bar--cobar counit. At all genera it is conditional on the
Francis--Gaitsgory $(\infty,2)$ descent package for the completed
ambient.

Therefore the BBL comparison is an equivalence whenever the counit
criterion holds and the required totalization preservation is
available. The reverse implication is not used as an independent
criterion.

\subsection{Kac--Shapovalov determinant}

\subsubsection*{\textup{(i)}~$\Rightarrow$~\textup{(ix)}}

If $\cA$ is chirally Koszul, the PBW filtration is strict on the
vacuum module: $\gr^F M_0(\cA) \cong \Symch(V)$ where
$V = F_1\cA/F_0\cA$. Strictness of the PBW map implies that the
Shapovalov form $G_h$ is non-degenerate in the bar-relevant range.

\subsubsection*{\textup{(ix)}~$\Rightarrow$~\textup{(i)}}

Conversely, non-degeneracy of $G_h$ in the bar-relevant range
means the PBW-to-bar comparison map
$\iota\colon\Symch(V)_h\to\barB_h(\cA)$ is injective at every
bar-relevant weight~$h$. Suppose for contradiction $d_r \neq 0$
for some $r \ge 2$. Then there exists
$x\in E_r^{p,q}$ with $d_r(x)\neq 0$; let $\tilde x\in F^p\barB$
be a lift. If $\tilde x\in\mathrm{im}(\iota)$, it represents a
nonzero class in $\gr^p(\barB)\cong \Sym^p(V)$, and
$d_r(\tilde x)$ represents a nonzero class in $\gr^{p+r}$. By
PBW strictness (injectivity of~$\iota$), this contradicts the
acyclicity of $\Sym(V)$ as a Koszul complex. Hence $d_r = 0$ for
all $r\ge 2$, giving (ii), hence~(i).

\subsection{The Lagrangian criterion, conditionally}

\subsubsection*{\textup{(i)}~$\Rightarrow$~\textup{(xi)}}

Under the perfectness and non-degeneracy hypotheses on the
ambient tangent complex~\cite[Prop.~6.3]{LorgatMKD1}: on the
Koszul locus, the bar-cobar adjunction provides a free resolution
of $\cA$, and the complementarity splitting
\[
\mathbf{Q}_g(\cA) \oplus \mathbf{Q}_g(\cA^!)
 \;\simeq\; H^*(\overline{\cM}_g, \cZ_\cA)
\]
identifies $\cM_\cA$ and $\cM_{\cA^!}$ as complementary subspaces
of $\cM_{\mathrm{comp}}$. The $(-1)$-shifted symplectic structure
on $\cM_{\mathrm{comp}}$ makes each of these subspaces isotropic;
complementarity gives maximal dimension, hence both are
Lagrangian.

\subsubsection*{\textup{(xi)}~$\Rightarrow$~\textup{(i)}}

If $\cM_\cA$ and $\cM_{\cA^!}$ are transverse Lagrangians and the
perfect comparison identifies their normal intersection complex
with the Koszul cone, then
their derived intersection
\[
\cM_\cA \times^h_{\cM_{\mathrm{comp}}} \cM_{\cA^!}
\]
is discrete: transverse Lagrangian intersection in a
$(-1)$-shifted symplectic space has expected dimension~$0$. The
normal comparison identifies discreteness of this intersection with
acyclicity of the chiral Koszul cone for the twisting datum
$(\cA,\barBgeom(\cA),\tau_0)$. That acyclicity is
Definition~\ref{def:chiral-koszul-morphism}.

\subsection{\textup{$\cD$}-module purity, one direction}

\subsubsection*{\textup{(xii)}~$\Rightarrow$~\textup{(x)}}

Suppose each $\barBgeom_n(\cA)$ is pure of weight~$n$ as a mixed
Hodge module. By Saito's theory~\cite{Saito}, the weight
filtration on a pure Hodge module splits up to semisimplification;
for the bar complex, this splitting implies that the restriction
$i_S^!\,\barB_n(\cA)$ to any FM boundary stratum $S$ is again
pure. Purity plus the alignment of the characteristic variety to
FM strata forces acyclicity outside degree~$0$ on every stratum,
which is condition~(x).

\begin{remark}[Converse direction]
\label{rem:d-module-purity-content}
The converse, that Koszulness implies $\cD$-module purity of the
bar complex, is established for the affine Kac--Moody family in
\cite[Prop.~7.5]{LorgatMKD1}: there the bar complex is identified
with a geometric complex whose purity is known via the Riemann
hypothesis for weight filtrations. For Virasoro and the $\cW$
algebras, the required geometric origin for the bar complex is not
known, and the converse is open.
\end{remark}

% ================================================================
\section{Examples}
\label{sec:examples}
% ================================================================

The following standard families give the first tests of the
intrinsic core.

\subsection{Heisenberg}

The Heisenberg vertex algebra $\Heis_k$ at level $k$ is generated
by a single current $J(z)$ with OPE
\[
J(z)J(w) \;\sim\; \frac{k}{(z-w)^2}.
\]
At nonzero level, the PBW filtration is strict: the associated
graded is a free polynomial Poisson vertex algebra. The
Kac--Shapovalov determinant is a product of factors
$(n - h)^{\mu_n}$ for the weight lattice, and is nonzero
generically. Condition~(ix) holds, and
Theorem~\ref{thm:koszul-equivalences-meta} supplies the other
core conditions on this row.

In the shadow tower language, $\Heis_k$ is a Gaussian algebra
(depth $r_{\max} = 2$): the bar complex is abelian and the
chiral Hochschild cohomology is concentrated in the single bar
degree $\{0,1,2\}$ that is visible on any free theory. The
$A_\infty$-structure is strictly associative, and the bar--cobar
counit is a quasi-isomorphism by direct computation.

\begin{remark}
The Heisenberg collision residue is
$r_{\Heis}^{\mathrm{coll}}(z)=k/z$. The shadow constant is
$\kappa^{\Heis}=k$. At $k=0$ the double-pole residue and
$\kappa$ vanish; the reduced bar remains the Gaussian free-field
bar complex.
\end{remark}

\subsection{Affine Kac--Moody at generic level}

Let $\hat\fg_k$ be the affine Kac--Moody vertex algebra of a
simple Lie algebra~$\fg$ at level~$k$. The Kac--Shapovalov
determinant is given by the Kac determinant formula; it vanishes
precisely at the Shapovalov walls $k + h^\vee = p/q$ for a
discrete set of rationals depending on the weight. Away from
these walls, condition~(ix) holds, and
Theorem~\ref{thm:koszul-equivalences-meta} gives chiral
Koszulness.

For the bar--cobar pair $(\hat\fg_k,\hat\fg_{k'})$ with
$k+k' = -2h^\vee$ (the Killing self-duality), the shadow
constant~$\kappa$ is
\[
\kappa^{\KM}(\fg,k)
 \;=\;\frac{\dim(\fg)\,(k+h^\vee)}{2h^\vee}.
\]
At the critical level $k = -h^\vee$, the shadow constant
$\kappa^{\KM}$ vanishes; the Shapovalov determinant acquires zeros
of the corresponding order, and the Kac--Shapovalov input
(condition (ix)) enters its degenerate regime. The critical level
lies outside the generic Koszul locus treated here. At generic
$k\neq -h^\vee$, $\kappa^{\KM}\neq 0$ and the PBW criterion
applies.

For $\fg = \mathfrak{sl}_2$ and $k = 1$,
\[
\kappa^{\KM}
 = \frac{3(1+2)}{2\cdot 2}
 = \frac{9}{4}.
\]
This agrees with the direct computation of the second shadow
coefficient from the OPE~\cite[\S3.1]{LorgatMKD1}, and at
$k = -h^\vee = -2$ one recovers $\kappa^{\KM} = 0$, reproducing
the critical-level degeneration. At level $k=0$ the trace-form
double-pole collision residue $k\Omega_{\mathrm{tr}}/z$ vanishes,
but the simple-pole Lie bracket channel remains. Thus
$V_0(\fg)$ is not Gaussian or abelian, and
$\kappa(V_0(\fg))=\dim(\fg)/2$.

\subsection{Virasoro in the Koszul locus}

The Virasoro vertex algebra $\Vir_c$ at central charge~$c$ has
shadow constant $\kappa^{\Vir} = c/2$. The Kac determinant formula
for $\Vir_c$ gives the discrete set of degenerate central charges
$c_{p,q} = 1 - 6(p-q)^2/(pq)$; away from these, the vacuum module
is irreducible in the bar-relevant range, and condition~(ix)
holds.

At $c = 13$, the Koszul pair $(\Vir_c,\Vir_{26-c})$ is self-dual,
and the duality constant $c + c' = 26$ (the Koszul conductor of
Virasoro) matches the matter plus $bc$ ghost anomaly-cancellation
normalization.
At these generic values, $\Vir_c$ is chirally Koszul and the
PBW/Kac criterion applies.

Virasoro has one strong generator. Hence
$\delta F_g^{\mathrm{cross}}(\Vir_c)=0$ for every~$g$, while the
class-M shadow tower remains infinite through the nonzero higher
Virasoro shadows.

\subsection{The \texorpdfstring{$\beta\gamma$}{betagamma} system}

The $\beta\gamma$ system is a free-field vertex algebra with
commuting fields $\beta(z), \gamma(z)$ satisfying
$\beta(z)\gamma(w)\sim 1/(z-w)$. Its associated graded is the
Poisson vertex algebra on $T^*(\cO_X)$ with the canonical
bracket; the bar complex is an exterior algebra on cotangent
generators.

The PBW bar model of the free $\beta\gamma$ system is formal:
the transferred higher $A_\infty$ operations vanish on the free
bar generators. This proves condition~\textup{(iii)}. The
shadow-tower statement is separate: $\beta\gamma$ is contact class
C, with depth $r_{\max}=4$ and a nonzero quartic contact shadow.

\subsection{Summary}

\begin{center}
\begin{tabular}{lll}
\toprule
Family & Decisive criterion & Conclusion \\
\midrule
Heisenberg & free PBW bar & Gaussian Koszul locus \\
Affine Kac--Moody & Kac--Shapovalov nondegeneracy & generic Koszul locus \\
Virasoro & PBW collapse off Kac walls & class-M Koszul locus \\
$\beta\gamma$ system & free $A_\infty$ bar formality & contact Koszul locus \\
\bottomrule
\end{tabular}
\end{center}

For these families, one decisive criterion enters the intrinsic
core. The Lagrangian criterion~\textup{(xi)} is available where
the perfectness and non-degeneracy hypotheses are verified.
$\cD$-module purity~\textup{(xii)} is proved for affine
Kac--Moody at non-critical level and remains open for Virasoro and
$\cW$.

% ================================================================
\section{Open questions}
\label{sec:openquestions}
% ================================================================

\subsection{The converse of $\cD$-module purity}

The remaining open equivalence is~(xii). The implication
$(\mathrm{xii})\Rightarrow(\mathrm{x})$ is unconditional
(Section~\ref{sec:proof}); the converse is known for the affine
Kac--Moody family (Remark~\ref{rem:d-module-purity-content}) but
open for Virasoro and $\cW$.

\begin{conjecture}
\label{conj:purity-converse}
For every chirally Koszul chiral algebra~$\cA$ on a smooth
projective curve, each $\barBgeom_n(\cA)$ is pure of weight~$n$
as a mixed Hodge module, with characteristic variety aligned to
the FM boundary strata.
\end{conjecture}

The obstacle is the absence of a known geometric model from which
purity of the Virasoro bar complex can be read off.  Any proof of
the converse must proceed either by a direct
computation of the weight filtration on $\barBgeom_n(\Vir_c)$, or
by an indirect route via a geometric model (e.g., a resolution of
the Virasoro bar complex by a pure complex on a related moduli
space).

\subsection{The perfectness hypothesis for (xi)}

Condition~(xi) requires perfectness and non-degeneracy of the
ambient tangent complex, which is verified for the standard
landscape~\cite[Prop.~6.3]{LorgatMKD1} but not known in full
generality. A natural question is whether perfectness holds for
every chirally Koszul algebra.

\begin{conjecture}
\label{conj:perfectness}
If $\cA$ is chirally Koszul, the ambient tangent complex at~$\cA$
in $\cM_{\mathrm{comp}}$ is perfect of tor-amplitude~$[0,2]$.
\end{conjecture}

If this conjecture holds, the Lagrangian criterion~(xi) is
equivalent to the intrinsic Koszulness core.

\subsection{The Koszul locus beyond the standard landscape}

A natural family of chiral algebras outside the standard landscape
consists of the non-principal $\cW$-algebras of type ADE. The
Koszul duality for hook-type $\cW$-algebras of type~$A$ has been
established in~\cite{LorgatMKD1}; other nilpotent orbits remain
open and are controlled by the associated variety hierarchy
recorded in~\cite{LorgatMKD1}.

\subsection{Shadow depth and Koszulness}

The shadow depth classification (Gaussian, Lie, contact, mixed)
and chiral Koszulness are logically distinct: Koszulness concerns
the diagonal concentration of the bar cohomology, while shadow
depth concerns the degree at which the shadow tower terminates.
The standard generic rows are chirally Koszul. Gaussian
(Heisenberg) and contact ($\beta\gamma$) rows have finite shadow
depth; Virasoro and $\cW$-algebras have $r_{\max}=\infty$
(class~M). The Koszulness criteria do not distinguish these
shadow classes.

\begin{remark}[Koszulness is not SC formality]
Chiral Koszulness is an assertion about $H^*$ in diagonal
bidegree. Swiss-cheese formality is the vanishing of higher
operations on the open-closed bar complex. The standard generic
families lie on the chiral Koszul locus; Swiss-cheese formality is
restricted to class~$G$.
\end{remark}

% ================================================================
\appendix
\section{Scope of the Auxiliary Criteria}
\label{sec:n1-scope}
% ================================================================

At genus~$0$, PBW $E_2$ collapse is the common hub for the Koszul
twisting datum, the bar--cobar counit, twisted tensor acyclicity,
Priddy bar-weight concentration, and the Hochschild amplitude
consequence. Swiss-cheese formality coincides with this hub only
on class~$G$; classes $L$, $C$, and $M$ may be Koszul and fail
Swiss-cheese formality.

The Barr--Beck--Lurie comparison is used as a monadic consequence
of the bar--cobar counit. Its all-genera form inherits the
Francis--Gaitsgory $(\infty,2)$ descent hypotheses of the completed
ambient. Factorization homology is a criterion only when the
filtered comparison $\mathsf{FH}$ and its detecting strengthening
$\mathsf{FH}^{\mathrm{det}}$ are available. Under the
uniform-weight hypothesis the surviving scalar class is
$\mathrm{obs}_g(\cA)=\kappa(\cA)\lambda_g$; multi-weight families
such as $\cW_N$ require the cross-channel correction beginning with
$\delta F_g^{\mathrm{cross}}$.

The Hochschild condition is one-way. PBW Koszulness gives
cohomological amplitude $[0,2]$ and duality for
$\ChirHoch^\bullet(\cA)$ on the Koszul locus. A free-polynomial
cup-product algebra is stronger and requires Massey vanishing for
the transferred Fulton--MacPherson product.

The Lagrangian criterion is equivalent to the intrinsic core only
when the ambient tangent complex is perfect and non-degenerate and
the normal-intersection comparison identifies the derived
intersection with the chiral Koszul cone. The $\cD$-module purity
criterion is proved in the direction
\textup{(xii)}$\Rightarrow$\textup{(x)}; the converse is known for
affine Kac--Moody at non-critical level and remains open in the
Virasoro and $\cW$ cases.

For admissible levels $k+h^\vee=p/q$, the periodic-CDG comparison
hypothesis supplies a Koszul-compatible filtration on the
non-degenerate simply-laced admissible case at the level of
cohomology of filtration quotients.
Strict chain-level periodicity still depends on Massey vanishing,
and descent from the small-quantum-group stabilization to the
semisimplified Kazhdan--Lusztig target remains a separate
conjectural step.

% ================================================================
\begin{thebibliography}{99}

\bibitem{BD04}
A.~Beilinson and V.~Drinfeld.
\newblock \emph{Chiral Algebras}.
\newblock American Mathematical Society Colloquium Publications,
 vol.~51, 2004.

\bibitem{BGS}
A.~Beilinson, V.~Ginzburg, and W.~Soergel.
\newblock Koszul duality patterns in representation theory.
\newblock \emph{J.~Amer.~Math.~Soc.} 9 (1996), 473--527.

\bibitem{CG}
K.~Costello and O.~Gwilliam.
\newblock \emph{Factorization algebras in quantum field theory},
 vols.~I--II.
\newblock Cambridge University Press, 2017/2021.

\bibitem{FG12}
J.~Francis and D.~Gaitsgory.
\newblock Chiral Koszul duality.
\newblock \emph{Selecta Math.~(N.S.)} 18 (2012), 27--87.

\bibitem{GR17}
D.~Gaitsgory and N.~Rozenblyum.
\newblock \emph{A study in derived algebraic geometry}.
\newblock American Mathematical Society, 2017.

\bibitem{HA}
J.~Lurie.
\newblock \emph{Higher Algebra}.
\newblock 2017. Available at
 \url{https://www.math.ias.edu/~lurie/}.

\bibitem{KacBook}
V.~Kac.
\newblock \emph{Infinite dimensional Lie algebras}.
\newblock Third edition, Cambridge University Press, 1990.

\bibitem{Keller}
B.~Keller.
\newblock Introduction to $A$-infinity algebras and modules.
\newblock \emph{Homology Homotopy Appl.} 3 (2001), 1--35.

\bibitem{Li}
H.~Li.
\newblock Abelianizing vertex algebras.
\newblock \emph{Comm.~Math.~Phys.} 259 (2005), 391--411.

\bibitem{LorgatMKD1}
R.~Lorgat.
\newblock \emph{Modular Koszul Duality, Volume~I}.
\newblock Monograph, 2026.

\bibitem{LV12}
J.-L.~Loday and B.~Vallette.
\newblock \emph{Algebraic Operads}.
\newblock Grundlehren der mathematischen Wissenschaften,
 vol.~346, Springer, 2012.

\bibitem{PP}
A.~Polishchuk and L.~Positselski.
\newblock \emph{Quadratic algebras}.
\newblock American Mathematical Society, 2005.

\bibitem{Priddy}
S.~Priddy.
\newblock Koszul resolutions.
\newblock \emph{Trans.~Amer.~Math.~Soc.} 152 (1970), 39--60.

\bibitem{Saito}
M.~Saito.
\newblock Mixed Hodge modules.
\newblock \emph{Publ.~Res.~Inst.~Math.~Sci.} 26 (1990), 221--333.

\end{thebibliography}

\end{document}
